\begin{helper}
Let $I$ be a nonempty finite set and let $s:I\to\mathbb R$ be any real-valued
function.  Then there is an element $i_{\max}\in I$ such that
$s(i)\le s(i_{\max})$ for every $i\in I$.
\end{helper}
\begin{proof}
Because $I$ is finite, the image $s(I)\subseteq\mathbb R$ is a finite subset
of the line.  Because $I$ is nonempty, this image is also nonempty.  A
nonempty finite subset of the line has a largest element: start with any
image value and, one by one, replace the current candidate by the next image
value if the next value is larger.  After finitely many steps the resulting
candidate is at least every element of $s(I)$.

Choose $i_{\max}\in I$ whose image is this largest value of $s(I)$.  For an
arbitrary $i\in I$, the number $s(i)$ belongs to the image $s(I)$, so by
maximality of $s(i_{\max})$ inside that finite image we have
$s(i)\le s(i_{\max})$.  This proves the claim.
\end{proof}
